% Paragraph 1 : Issue
Estimating the governing parameters of nonlinear dynamical systems from time-series observations is a fundamental inverse problem with broad applications in epidemiology, ecology, and physiology. % TODO: citation (fundamental inverse problems) 
In many practical settings, however, observations are severely constrained: they are often \emph{partial} (only a subset of the state variables is measurable) and \emph{sparse} (a relatively small number of time points are observed). 
These physical constraints render the inverse mapping from observations to parameters highly ill-posed and, in certain regimes, structurally non-identifiable. 
When key states remain hidden, resolving these unobserved latent trajectories simultaneously with the system parameters formulates a joint optimization problem.
This coupling inherently creates a highly non-convex loss landscape, drastically amplifying the computational and theoretical difficulty.

% Paragraph 2 : Limitations of existing methods
Traditional approaches combine mechanistic numerical integration with optimization techniques, such as nonlinear least squares (NLLS) or sampling based method (MCMC). % TODO: citation (NLLS/MCMC textbooks)
While mathematically principled, these per-instance optimizations incur significant computational costs and are highly sensitive to initial conditions; 
they frequently fall into local minima or diverge when applied to sparse measurements of complex oscillatory or non-autonomous systems.

% Paragraph 3 : Propose our solution
To address these limitations, we propose a novel parameter estimation framework structured as a contractive fixed-point iteration. 
We first decouple the joint inverse problem into two conditional tasks: a hidden-state estimator infers the latent trajectory given partial observations and a parameter guess, while a parameter estimator updates the parameter vector given the observations and the inferred hidden state. 
Inference is then performed by alternating updates between these two trained neural networks. However, simply coupling two neural networks offers no theoretical guarantee of stability; the alternating sequence could easily oscillate or diverge. 
To overcome this, we enforce Lipschitz bounds on both networks via spectral normalization. 
This rigorous structural constraint guarantees that their composition forms a contraction mapping, ensuring that the inference iteration converges exponentially to a unique fixed point from any initialization.

% Paragraph 4 : Results
We demonstrate the effectiveness of this framework through extensive numerical evaluations on epidemiological (SIR), ecological (Lotka-Volterra), and physiological (OGTT) models. 
Our results show that the proposed fixed-point architecture robustly recovers parameters from sparse and partial data in regimes where traditional joint optimization fails. 
To provide a deeper understanding of these empirical outcomes, we utilize the Hermann-Krener nonlinear observability framework~\cite{hermann2003nonlinear} as an analytical tool. 
This theoretical grounding allows us to explicitly characterize identifiable regimes and provides a principled explanation for phenomena such as regression-to-mean (model collapse) encountered in structurally unidentifiable scenarios.

%\textbf{Contributions.} The main contributions of this paper are threefold:
%\begin{itemize}
%    \item We introduce a decoupled fixed-point iteration architecture that enables stable parameter inference from sparse and partial observations.
%    \item We provide a rigorous mathematical guarantee that the proposed iterative update acts as a contraction mapping, ensuring exponential convergence to a unique fixed point.
%    \item We empirically validate the framework on multiple nonlinear systems and provide a theoretical analysis of structural identifiability to clarify the boundaries of reliable parameter recovery.
%\end{itemize}